\chapter{自定义组件}
\label{chap:self-defined}

介绍如何配置和使用一些自定义组件。

\section{公式高亮}

\vspace{3em}
\begin{equation*}
    Attention(\eqnmarkbox[green]{Q}{Q}, \eqnmarkbox[blue]{K}{K}, \eqnmarkbox[red]{V}{V}) = softmax(\frac{QK^T}{\eqnmarkbox[purple]{sqrt_dk}{\sqrt{d_k}}})V
\end{equation*}
\annotate[yshift=3em]{above}{Q}{Query Matrix}
\annotate[yshift=1em]{above}{K}{Key Matrix}
\annotate[yshift=-1em]{below}{V}{Value Matrix}
\annotate[yshift=-0.1em]{below}{sqrt_dk}{The Normalizer}
\vspace{1em}

\begin{lstlisting}[language=tex, caption={公式高亮示例代码}]
\vspace{4em}
\begin{equation*}
    Attention(\eqnmarkbox[green]{Q}{Q}, \eqnmarkbox[blue]{K}{K}, \eqnmarkbox[red]{V}{V}) = softmax(\frac{QK^T}{\eqnmarkbox[purple]{sqrt_dk}{\sqrt{d_k}}})V
\end{equation*}
\annotate[yshift=3em]{above}{Q}{Query Matrix}
\annotate[yshift=1em]{above}{K}{Key Matrix}
\annotate[yshift=-1em]{below}{V}{Value Matrix}
\annotate[yshift=-0.1em]{below}{sqrt_dk}{The Normalizer}
\vspace{1em}
\end{lstlisting}

\section{彩色定理框}
\index{定理框}

本模板使用 \texttt{tcolorbox} 宏包提供了多种彩色定理框环境。

\subsection{定义}

\begin{definition}{欧拉公式}{euler}
    对于任意实数 $x$，有
    \[ e^{ix} = \cos x + i \sin x \]
    特别地，当 $x = \pi$ 时，得到著名的欧拉恒等式：$e^{i\pi} + 1 = 0$。
\end{definition}

引用定义：\cref{def:euler}。

\subsection{定理与引理}

\begin{lemma}{三角不等式}{triangle}
    对于任意复数 $z_1, z_2$，有
    \[ |z_1 + z_2| \leq |z_1| + |z_2| \]
\end{lemma}

\begin{theorem}{勾股定理}{pythagorean}
    设直角三角形的两条直角边长分别为 $a$ 和 $b$，斜边长为 $c$，则
    \[ a^2 + b^2 = c^2 \]
\end{theorem}

引用引理和定理：\cref{lem:triangle}，\cref{thm:pythagorean}。

\subsection{推论与例}

\begin{corollary}{单位圆上的点}{unit-circle}
    由欧拉公式可知，$e^{i\theta}$ 表示单位圆上幅角为 $\theta$ 的点。
\end{corollary}

\begin{example}{计算三角形斜边}{calc-hypotenuse}
    已知直角三角形两条直角边长分别为 3 和 4，求斜边长。

    \textbf{解}：根据勾股定理，$c = \sqrt{3^2 + 4^2} = \sqrt{9 + 16} = \sqrt{25} = 5$。
\end{example}

\subsection{注}

\begin{remark}
    勾股定理是欧几里得几何中最基本的定理之一，在中国古代称为"商高定理"或"勾三股四弦五"。
\end{remark}

\section{提示框}
\index{提示框}

本模板提供了三种提示框，用于突出显示重要信息。

\begin{tipbox}
    编译本模板推荐使用 \texttt{latexmk} 命令，它会自动处理多次编译和 biber 调用。
\end{tipbox}

\begin{warnbox}
    请确保使用 XeLaTeX 引擎编译本文档，pdfLaTeX 可能无法正确处理中文。
\end{warnbox}

\begin{notebox}
    如果遇到参考文献不显示的问题，请检查是否正确运行了 biber 命令。
\end{notebox}

\section{算法伪代码}
\index{算法}

本模板使用 \texttt{algorithm2e} 宏包编写算法，支持中文关键字。

\begin{algorithm}[H]
    \caption{二分查找}
    \label{alg:binary-search}
    \KwIn{有序数组 $A[0..n-1]$，目标值 $x$}
    \KwOut{目标值的索引，若不存在则返回 $-1$}
    $l \leftarrow 0$\;
    $r \leftarrow n - 1$\;
    \While{$l \leq r$}{
        $m \leftarrow \lfloor (l + r) / 2 \rfloor$\;
        \If{$A[m] = x$}{
            \KwRet{$m$}\;
        }
        \ElseIf{$A[m] < x$}{
            $l \leftarrow m + 1$\;
        }
        \Else{
            $r \leftarrow m - 1$\;
        }
    }
    \KwRet{$-1$}\;
\end{algorithm}

\vspace{1em}

\begin{algorithm}[H]
    \caption{快速排序}
    \label{alg:quicksort}
    \KwIn{数组 $A[p..r]$}
    \KwOut{排序后的数组}
    \If{$p < r$}{
        $q \leftarrow$ \textsc{Partition}$(A, p, r)$\;
        \textsc{QuickSort}$(A, p, q-1)$\;
        \textsc{QuickSort}$(A, q+1, r)$\;
    }
\end{algorithm}
